\documentclass[tikz,border=4pt]{standalone}
\usepackage{amsmath}
\usepackage{sfmath}
\usetikzlibrary{arrows.meta,calc,decorations.pathmorphing,patterns,positioning}

\definecolor{sulfur}{HTML}{E5B90B}
\definecolor{sulfurdeep}{HTML}{9C7500}
\definecolor{polymer}{HTML}{343A40}
\definecolor{charge}{HTML}{C53A4A}
\definecolor{fieldblue}{HTML}{2878B5}
\definecolor{trapblue}{HTML}{305A8A}
\definecolor{softblue}{HTML}{EAF2F8}
\definecolor{softgold}{HTML}{FFF5CF}
\definecolor{softred}{HTML}{FBEAEC}
\definecolor{apparatus}{HTML}{66727D}
\definecolor{ink}{HTML}{20262C}
\definecolor{muted}{HTML}{65717C}

\tikzset{
  >=Latex,
  every node/.style={font=\sffamily, text=ink},
  bond/.style={draw=polymer, line width=0.85pt, line cap=round},
  sulfur bead/.style={circle, draw=sulfurdeep, fill=sulfur, line width=0.55pt,
    minimum size=3.2mm, inner sep=0pt, font=\sffamily\scriptsize\bfseries},
  charge bead/.style={circle, draw=charge, fill=white, text=charge,
    line width=0.65pt, minimum size=3.6mm, inner sep=0pt,
    font=\sffamily\scriptsize\bfseries},
  axis/.style={draw=ink, line width=0.65pt, -Latex},
  trace/.style={draw=fieldblue, line width=1.25pt, line cap=round},
  apparatus line/.style={draw=apparatus, line width=0.85pt, line cap=round},
  callout/.style={rounded corners=1.2pt, draw=muted!55, fill=white,
    inner sep=2.5pt, font=\sffamily\scriptsize, align=center},
  panel title/.style={font=\sffamily\small\bfseries, anchor=west},
  tiny label/.style={font=\sffamily\scriptsize, text=muted, align=center},
}

\newcommand{\panel}[3]{%
  \draw[rounded corners=2.2pt, draw=ink!32, fill=white, line width=0.55pt]
    (0,0) rectangle (8.55,5.40);
  \node[font=\sffamily\large\bfseries, anchor=north west] at (0.18,5.25) {#1};
  \node[panel title] at (0.72,5.05) {#2};
  \node[tiny label, anchor=east] at (8.28,5.06) {#3};
}

\begin{document}
\begin{tikzpicture}[x=1cm,y=1cm]

% -----------------------------------------------------------------------------
% A. Primary microstructure
% -----------------------------------------------------------------------------
\begin{scope}[shift={(0,5.75)}]
  \panel{A}{poly(S-r-DIB) primary microstructure}{inverse-vulcanized network motif}

  \node[tiny label, anchor=west] at (0.45,4.48) {sulfur-rich segment};
  \foreach \x in {1.00,1.50,2.00,2.50,3.00} {
    \node[sulfur bead] at (\x,3.92) {S};
  }
  \draw[bond] (0.55,3.92)--(0.84,3.92);
  \draw[bond] (1.16,3.92)--(1.34,3.92);
  \draw[bond] (1.66,3.92)--(1.84,3.92);
  \draw[bond] (2.16,3.92)--(2.34,3.92);
  \draw[bond] (2.66,3.92)--(2.84,3.92);
  \draw[bond] (3.16,3.92)--(3.52,3.92);
  \node[font=\sffamily\scriptsize, text=muted] at (2.00,3.48) {$\mathrm{S}_{n}$};

  % Stylized DIB aromatic crosslinker
  \coordinate (c) at (5.10,3.88);
  \foreach \a in {0,60,...,300} {
    \coordinate (p\a) at ($(c)+(\a:0.72)$);
  }
  \draw[bond, line width=1pt]
    (p0)--(p60)--(p120)--(p180)--(p240)--(p300)--cycle;
  \draw[bond, line width=0.55pt]
    ($(p0)!0.22!(p60)$)--($(p0)!0.78!(p60)$)
    ($(p120)!0.22!(p180)$)--($(p120)!0.78!(p180)$)
    ($(p240)!0.22!(p300)$)--($(p240)!0.78!(p300)$);
  \draw[bond] (p120)--++(150:0.58)--++(90:0.40);
  \draw[bond] (p240)--++(210:0.58)--++(-90:0.40);
  \draw[bond] (p0)--++(0:0.58)--++(90:0.40);
  \node[font=\sffamily\scriptsize\bfseries] at (5.10,2.70) {DIB-derived junction};

  \draw[bond] (3.52,3.92) .. controls (3.90,4.32) and (4.05,4.34) .. (p180);
  \draw[bond] (p0) .. controls (6.28,4.30) and (6.45,4.25) .. (6.88,4.00);
  \foreach \x in {7.02,7.52} {\node[sulfur bead] at (\x,4.00) {S};}
  \draw[bond] (6.88,4.00)--(6.86,4.00);
  \draw[bond] (7.18,4.00)--(7.36,4.00);
  \draw[bond] (7.68,4.00)--(8.00,4.00);

  \draw[decorate, decoration={brace, amplitude=4pt}, draw=muted]
    (0.72,2.02)--(7.88,2.02)
    node[midway, below=6pt, font=\sffamily\scriptsize, text=muted]
    {random sulfur--organic connectivity};
  \node[callout, fill=softgold] at (4.30,1.15)
    {high sulfur content supplies polarizable, structurally diverse\charge-localization environments};
\end{scope}

% -----------------------------------------------------------------------------
% B. Composition series
% -----------------------------------------------------------------------------
\begin{scope}[shift={(8.90,5.75)}]
  \panel{B}{Sulfur-composition variation}{representative chains}

  \foreach \yy/\lab/\count/\shade in {4.13/S60/4/25,3.03/S75/6/55,1.93/S85/8/85} {
    \node[font=\sffamily\small\bfseries, anchor=east] at (1.05,\yy) {\lab};
    \draw[bond] (1.35,\yy)--(1.62,\yy);
    \foreach \i in {1,...,\count} {
      \pgfmathsetmacro{\xx}{1.45+0.55*\i}
      \node[sulfur bead, fill=sulfur!\shade] at (\xx,\yy) {S};
      \ifnum\i>1
        \pgfmathsetmacro{\xp}{1.45+0.55*(\i-1)}
        \draw[bond] (\xp+0.17,\yy)--(\xx-0.17,\yy);
      \fi
    }
    \pgfmathsetmacro{\endx}{1.45+0.55*\count}
    \draw[bond] (\endx+0.17,\yy)--(7.35,\yy);
    \draw[bond] (7.35,\yy)--(7.56,\yy+0.18)--(7.77,\yy)--(7.98,\yy+0.18);
  }

  \draw[-Latex, draw=sulfurdeep, line width=1.1pt] (2.05,1.10)--(7.55,1.10);
  \node[font=\sffamily\scriptsize\bfseries, text=sulfurdeep] at (4.80,1.42)
    {increasing sulfur fraction};
  \node[font=\sffamily\scriptsize, text=muted] at (4.80,0.61)
    {series: S60 / S70 / S75 / S80 / S85};
\end{scope}

% -----------------------------------------------------------------------------
% C. Localized traps: real and energy space
% -----------------------------------------------------------------------------
\begin{scope}[shift={(0,0)}]
  \panel{C}{Localized charge traps}{real space $\leftrightarrow$ energy space}

  \node[font=\sffamily\scriptsize\bfseries] at (2.18,4.55) {real-space localization};
  \fill[softgold, rounded corners=3pt] (0.52,1.05) rectangle (4.02,4.25);
  \foreach \x/\y in {0.85/1.50,1.18/3.55,1.62/2.55,2.20/3.72,2.56/1.45,3.05/2.85,3.62/3.62} {
    \node[sulfur bead, minimum size=2.8mm] at (\x,\y) {S};
  }
  \draw[bond] (0.85,1.50)..controls(1.25,1.90)..(1.62,2.55)
    ..controls(1.85,3.05)..(1.18,3.55);
  \draw[bond] (2.20,3.72)..controls(2.48,3.18)..(3.05,2.85)
    ..controls(3.42,2.48)..(3.62,3.62);
  \draw[bond] (2.56,1.45)..controls(2.66,2.15)..(3.05,2.85);
  \node[charge bead] (q1) at (1.80,1.65) {$+$};
  \node[charge bead] (q2) at (2.72,3.23) {$+$};
  \draw[draw=charge!60, dashed, line width=0.7pt] (q1) circle[radius=0.36];
  \draw[draw=charge!60, dashed, line width=0.7pt] (q2) circle[radius=0.36];
  \node[tiny label] at (2.28,0.72) {spatially heterogeneous trap sites};

  \draw[axis] (4.72,1.00)--(4.72,4.33) node[above, font=\sffamily\scriptsize] {$E$};
  \draw[apparatus line] (4.72,3.65)--(8.00,3.65);
  \draw[apparatus line] (4.72,1.55)--(8.00,1.55);
  \node[font=\sffamily\scriptsize, anchor=west] at (4.84,3.84) {$E_{\mathrm{c}}$};
  \node[font=\sffamily\scriptsize, anchor=west] at (4.84,1.22) {$E_{\mathrm{v}}$};
  \foreach \x/\y in {5.35/3.08,6.00/2.72,6.55/2.33,7.20/2.95,7.63/2.12} {
    \draw[draw=trapblue, line width=1.05pt] (\x-0.17,\y)--(\x+0.17,\y);
    \draw[-Latex, draw=charge, line width=0.75pt] (\x,3.60)--(\x,\y+0.09);
  }
  \node[tiny label] at (6.38,0.72) {distributed trap depths $E_{\mathrm{t}}$};
\end{scope}

% -----------------------------------------------------------------------------
% D. Kinetic evidence: SMU/MIM and current decay
% -----------------------------------------------------------------------------
\begin{scope}[shift={(8.90,0)}]
  \panel{D}{Kinetic evidence}{SMU / MIM current decay}

  \node[font=\sffamily\scriptsize\bfseries] at (2.13,4.55) {metal--insulator--metal};
  \draw[rounded corners=2pt, draw=apparatus, fill=softblue, line width=0.75pt]
    (0.48,2.58) rectangle (1.70,3.64);
  \node[font=\sffamily\scriptsize\bfseries, align=center] at (1.09,3.11) {SMU};
  \draw[apparatus line] (1.70,3.35)-|(2.38,4.00);
  \draw[apparatus line] (1.70,2.83)-|(2.38,1.78);
  \fill[apparatus] (2.38,3.75) rectangle (3.82,4.02);
  \fill[sulfur!42] (2.38,2.03) rectangle (3.82,3.75);
  \fill[apparatus] (2.38,1.76) rectangle (3.82,2.03);
  \node[font=\sffamily\scriptsize] at (3.10,2.88) {polymer};
  \node[tiny label] at (3.10,1.38) {bias $\rightarrow$ polarization / detrapping};
  \draw[-Latex, draw=charge, line width=0.9pt] (2.80,3.40)--(2.80,2.42);
  \draw[-Latex, draw=charge, line width=0.9pt] (3.40,3.40)--(3.40,2.62);

  \draw[axis] (4.55,1.35)--(8.05,1.35) node[right, font=\sffamily\scriptsize] {$t$};
  \draw[axis] (4.55,1.35)--(4.55,4.35) node[above, font=\sffamily\scriptsize] {$I(t)$};
  \draw[trace] plot[smooth, domain=4.62:7.85, samples=80]
    (\x,{1.58+2.38*exp(-1.08*(\x-4.62))});
  \draw[draw=charge, dashed, line width=0.75pt] (4.62,1.58)--(7.90,1.58);
  \node[callout, anchor=north] at (6.32,4.18)
    {$I(t)=I_{\infty}+\displaystyle\sum_i A_i e^{-t/\tau_i}$};
  \node[tiny label] at (6.30,0.67) {relaxation times report trap kinetics};
\end{scope}

% -----------------------------------------------------------------------------
% E. ISPD: corona/probe, potential decay, trap-energy distribution
% -----------------------------------------------------------------------------
\begin{scope}[shift={(0,-5.75)}]
  \panel{E}{ISPD evidence}{potential decay $\rightarrow$ trap spectrum}

  \node[font=\sffamily\scriptsize\bfseries] at (1.72,4.57) {corona charge / probe};
  \fill[sulfur!35] (0.48,1.45) rectangle (3.45,1.86);
  \fill[apparatus!80] (0.48,1.25) rectangle (3.45,1.45);
  \draw[apparatus line, line width=1.05pt] (1.17,4.10)--(1.17,3.20);
  \draw[apparatus line] (0.94,4.10)--(1.40,4.10)--(1.17,4.46)--cycle;
  \foreach \x in {0.80,1.12,1.44} {
    \node[charge bead, minimum size=2.8mm] at (\x,2.52) {$+$};
    \draw[-Latex, draw=charge, line width=0.65pt] (\x,3.08)--(\x,2.70);
  }
  \draw[apparatus line, rounded corners=2pt] (2.35,3.62)--(3.07,3.62)--(3.07,3.08)--(2.35,3.08)--cycle;
  \draw[apparatus line] (2.71,3.08)--(2.71,2.24);
  \draw[apparatus line] (2.48,2.24)--(2.94,2.24);
  \node[tiny label] at (2.72,4.02) {non-contact probe};
  \node[tiny label] at (1.95,0.84) {surface-charged polymer};

  \draw[axis] (3.92,2.72)--(6.05,2.72) node[right, font=\sffamily\scriptsize] {$t$};
  \draw[axis] (3.92,2.72)--(3.92,4.54) node[above, font=\sffamily\scriptsize] {$V_{\!s}$};
  \draw[trace] plot[smooth, domain=3.98:5.92, samples=60]
    (\x,{2.90+1.38*exp(-1.35*(\x-3.98))});
  \node[tiny label] at (4.98,2.32) {surface-potential decay};

  \draw[axis] (6.38,1.20)--(8.14,1.20) node[right, font=\sffamily\scriptsize] {$E_{\mathrm{t}}$};
  \draw[axis] (6.38,1.20)--(6.38,4.54) node[above, font=\sffamily\scriptsize] {$N(E_{\mathrm{t}})$};
  \fill[trapblue!18] plot[smooth, domain=6.42:8.02, samples=60]
    (\x,{1.20+2.45*exp(-5.8*(\x-7.30)^2)}) -- (8.02,1.20)--(6.42,1.20)--cycle;
  \draw[draw=trapblue, line width=1.15pt] plot[smooth, domain=6.42:8.02, samples=60]
    (\x,{1.20+2.45*exp(-5.8*(\x-7.30)^2)});
  \draw[-Latex, draw=muted, line width=0.75pt] (5.72,1.83)--(6.28,1.83);
  \node[tiny label] at (7.28,0.68) {derived trap-energy distribution};
\end{scope}

% -----------------------------------------------------------------------------
% F. Floating cantilever and driven electrode
% -----------------------------------------------------------------------------
\begin{scope}[shift={(8.90,-5.75)}]
  \panel{F}{Trapped-charge Coulomb repulsion}{mechanically held, electrically floating}

  % Mechanical restraint only: deliberately no electrical connector.
  \fill[muted!18, rounded corners=1pt] (0.52,2.76) rectangle (1.46,4.30);
  \draw[draw=muted, line width=0.75pt] (0.52,2.76) rectangle (1.46,4.30);
  \foreach \y in {2.88,3.12,...,4.20} {
    \draw[draw=muted!55, line width=0.45pt] (0.57,\y)--(0.86,\y+0.29);
  }
  \draw[draw=ink, line width=1.2pt] (1.46,3.33)--(1.82,3.33);
  \node[tiny label, anchor=north] at (0.99,2.64) {mechanical\restraint};

  % Cantilever curves upward, away from the driven electrode.
  \path[draw=polymer, fill=sulfur!38, line width=1pt]
    (1.82,3.15) .. controls (3.72,3.10) and (5.15,3.66) .. (6.18,4.18)
    -- (6.18,4.48) .. controls (4.96,3.92) and (3.60,3.44) .. (1.82,3.45) -- cycle;
  \node[font=\sffamily\scriptsize\bfseries] at (4.10,3.67) {floating polymer cantilever};
  \foreach \x/\y in {2.38/3.46,3.15/3.53,3.92/3.70,4.70/3.96,5.45/4.27} {
    \node[charge bead, minimum size=3.0mm] at (\x,\y) {$+$};
  }

  % Driven electrode and a clearly visible air gap.
  \fill[apparatus] (1.85,1.48) rectangle (6.35,1.82);
  \node[font=\sffamily\scriptsize\bfseries, text=white] at (4.10,1.65) {driven electrode};
  \draw[<->, draw=muted, line width=0.65pt] (6.02,1.90)--(6.02,3.72);
  \node[tiny label, fill=white, inner sep=1pt] at (6.02,2.78) {air gap};
  \foreach \x in {2.38,3.15,3.92,4.70,5.45} {
    \node[charge bead, minimum size=3.0mm] at (\x,1.98) {$+$};
  }
  \draw[-{Latex[length=3.2mm]}, draw=charge, line width=1.35pt]
    (5.32,3.15)--(5.32,4.70)
    node[above, font=\sffamily\scriptsize\bfseries, text=charge] {$F_{\mathrm{C}}$};

  % Compact source circuit: electrode lead, source, and return terminating at ground.
  \draw[apparatus line] (6.35,1.65)-|(7.12,2.12);
  \draw[draw=apparatus, fill=softblue, line width=0.8pt, rounded corners=2pt]
    (6.72,0.72) rectangle (7.52,1.72);
  \node[font=\sffamily\scriptsize\bfseries] at (7.12,1.22) {$V$};
  \node[font=\sffamily\tiny, text=muted] at (7.12,0.91) {source};
  \draw[apparatus line] (7.12,0.72)--(7.12,0.46);
  \draw[apparatus line] (6.78,0.46)--(7.46,0.46);
  \draw[apparatus line] (6.88,0.31)--(7.36,0.31);
  \draw[apparatus line] (7.00,0.17)--(7.24,0.17);
  \node[tiny label, anchor=west] at (7.56,0.46) {source return};

  \node[callout, fill=softred, anchor=west] at (0.55,0.72)
    {trapped like charges repel;\the sample has no electrical ground connection};
\end{scope}

\end{tikzpicture}
\end{document}
